% Cover letter for Physica D: Nonlinear Phenomena. Compile with pdflatex.
%
% This letter goes to the editor and not to the referees, so it carries the
% author's name and the repository URL.
%
% Written after a desk rejection from Nonlinearity on the grounds that the
% manuscript did not make a sufficient contribution. The previous version of
% this letter described the contribution as "deliberately narrow", which handed
% the editor that conclusion in the second paragraph. Every attribution is
% retained below, because those are matters of fact and of credit. What is gone
% is the author's own verdict on significance, which is the editor's to reach.
\documentclass[11pt]{article}
\usepackage[margin=1.2in]{geometry}
\usepackage{parskip}
\pagestyle{empty}

\begin{document}

\noindent To the Editors of \emph{Physica D: Nonlinear Phenomena}

\bigskip

\noindent Re: ``The frozen blowup profile of the Okamoto, Sakajo and Wunsch
family: a control variate for its spectral exponent''

\bigskip

\noindent Dear Editors,

I am submitting the above manuscript for consideration as a research paper.

The paper concerns the Okamoto, Sakajo and Wunsch family of one dimensional
models for the Euler vorticity equation, and the frozen blowup profile it
admits on the circle above a critical advection. The Fourier spectrum of that
profile decays algebraically with an exponent $\beta$, which has until now been
obtained by least squares fits to the spectrum. I show that such fits are
accurate only to about a percent and are biased low, and give a determination
that fits nothing: $\beta$ follows from a single local quantity, read off the
computed profile by using a second constant, known exactly, as a control
variate to cancel the discretisation error the two share. This returns five
digits.

Three things follow. Spectral fitting is the standard route to this exponent
throughout the literature, so the bias attaches to published values: the paper
corrects one, giving $\beta = 9.29546$ at $a = 0.71$ against a published
$9.32592$. The corrected exponents vary continuously with $a$ and exclude the
integer value $3$, which is a statement about the structure of these profiles
rather than about precision, and which a fitted exponent is too coarse to
establish. And the error analysis gives a checkable condition, that the
discretisation error be of rank one, under which cancelling one constant
against another is legitimate at all; the test transfers wherever the technique
is tried.

Two disclosures. The setting is not mine: the frozen profile was found
numerically by Lushnikov, Silantyev and Siegel, and its existence, stability
and local exponent near $a = 1$ were proved by Chen, and the introduction and a
Limitations section state what is theirs and what is mine. Second, the paper
corrects a value published by Lushnikov, Silantyev and Siegel, and I wrote to
them describing the discrepancy and offering the manuscript before submitting
it anywhere.

Every number in the paper is reproducible. The solver, the profile solver, the
scripts behind each table and a suite of nineteen checks against closed forms
and conserved quantities are openly available under an MIT licence at
\texttt{github.com/munawarkazmi/degregorio-blowup}, and the checks run on every
commit.

The manuscript is not under consideration elsewhere, has not been published
previously, and there are no competing interests. The work was carried out
independently, with no institutional funding.

\bigskip

\noindent Yours faithfully,

\bigskip

\noindent Munawar Kazmi\\
\texttt{munawarkazmi.com}\\
ORCID 0009-0004-6565-6067

\end{document}
